% Appendix: Errata for the first draft of this report.
% \input this file after \appendix in main.tex.

\chapter{Errata: Changes from the First Draft}
\label{app:errata}

The first draft of this report (\emph{Lightweight Session-Based Recommendation
System}, YooChoose only) contained methodological issues that invalidate its
reported numbers, as well as several reporting errors. All experiments in this
report were re-run from scratch under the corrected protocol. This appendix
records what was wrong and how it was fixed, so that the earlier results are not
compared with the present ones.

\section{Methodological Issues (Results Re-run)}

\begin{enumerate}[leftmargin=*]
  \item \textbf{Data leakage through a sample-level split.}
        Each session was expanded into prefix samples
        $\{(i_1)\!\to\! i_2, (i_1,i_2)\!\to\! i_3,\ldots\}$ and the samples were
        then shuffled and split 80/20. Prefixes of the same session therefore
        appeared in both training and test sets: a test sample
        $(i_1,\ldots,i_4)\!\to\! i_5$ could coexist with the training sample
        $(i_1,\ldots,i_5)\!\to\! i_6$, whose input contains the test target.
        \emph{Fix:} split by session and by time, following the standard
        GRU4Rec protocol (sessions of the final day form the test set, the day
        before forms the validation set).

  \item \textbf{Test set used for model selection.}
        Per-epoch ``validation'' loss and metrics were computed on the test set,
        and the final reported row came from the same set.
        \emph{Fix:} separate train / validation / test splits; early stopping and
        hyperparameter choices use validation only; test metrics are computed
        once.

  \item \textbf{Non-standard sampling of YooChoose.}
        50{,}000 sessions were sampled at random from the full six-month log,
        which yields a sparse vocabulary and prevents comparison with published
        results. A figure caption nevertheless described the data as the
        ``YooChoose 1/64'' split. An item-frequency filter was mentioned but not
        applied.
        \emph{Fix:} the standard yoochoose1/64 (and 1/4) splits with item support
        $\geq 5$ and session length $\geq 2$.

  \item \textbf{Inconsistent dataset statistics.}
        A mean training prefix length of 4.85 implies sessions of roughly ten
        items, which would generate far more than the reported 128{,}029 samples
        from 50{,}000 sessions. The item count was also off by one: the model had
        an embedding table of 11{,}910 rows (11{,}909 items plus a padding row)
        and an output layer of 11{,}909 classes, so the dataset contained
        11{,}909 items rather than 11{,}910.
        \emph{Fix:} all statistics are now computed and exported by the notebook.

  \item \textbf{The ``cubic B-spline'' was a step function.}
        Executing the published B-spline routine (\texttt{\_bspline}) on a dense grid
        shows that every non-zero basis function is an indicator of a single
        grid interval (value 1 inside, 0 outside), i.e.\ a degree-0 spline, and
        that 3 of the 11 coefficients per edge are never used. The recursion
        reused one $\alpha$ for both terms and never extended the knot vector, so
        it did not implement the Cox--de~Boor formula given in the report.
        The learned spline functions were therefore piecewise constant, so the
        spline path had zero gradient with respect to its input; only the
        parallel linear path carried input gradients.
        \emph{Fix:} a correct B-spline basis on an extended grid, verified
        against a reference implementation (partition of unity and pointwise
        agreement), together with RBF and Chebyshev bases.

  \item \textbf{Ungated recurrent KAN cell.}
        The Pure KAN cell $h_t=\tanh(\mathrm{LN}(\mathrm{KAN}([x_t,h_{t-1}])))$
        has no update or reset gate, making it a vanilla RNN with a different
        non-linearity, and its input was squashed twice ($\tanh$ inside
        \texttt{KANLinear} and again after LayerNorm). It was also unrolled in a
        Python loop, which accounts for most of the reported 10.8$\times$
        slowdown.
        \emph{Fix:} KAN placements that keep GRU gating (Chapter~4), and the
        original Pure KAN and Hybrid models re-run under the corrected protocol
        for reference.

  \item \textbf{No control for added capacity.}
        KAN models had 25--43\% more parameters than the baseline, so any gain
        could not have been attributed to the KAN itself.
        \emph{Fix:} every KAN variant is paired with a parameter-matched MLP.
\end{enumerate}

\section{Reporting Errors}

\begin{enumerate}[leftmargin=*]
  \item The final comparison table listed Recall@20 with the HR@10 values
        (0.5808 / 0.5174 / 0.5386). With a single target, Recall@20 equals
        HR@20: 0.6542 / 0.5850 / 0.6071.
  \item The validation-loss deltas had the wrong sign. KAN and Hybrid had
        \emph{higher} loss than the RNN: $+0.2829$ and $+0.1525$.
  \item The table declared the RNN the winner on ``7/7 metrics'' while listing
        8 metrics.
  \item The Hybrid model's training time is 62.5\% of the Pure KAN time
        (981.6\,s / 1569.8\,s), not 67\%.
  \item Pure KAN first reached HR@10 $\geq 0.50$ at epoch 12 (0.5037), not
        epoch 14.
  \item Captions stated that KAN and Hybrid ``continue improving'' through
        epoch 20. The logs show the Hybrid validation loss reaching its minimum
        at epoch 12 (5.7161) and both models plateauing by epoch 16, as
        Chapter 5 of the draft itself stated.
  \item One figure was labelled Recall@10 while the tables reported Recall@20.
  \item The MRR@20 definition in the metrics table (average over sessions whose
        target is in the top 20) disagreed with the formula (average over all
        sessions, with misses counted as zero). The formula is correct.
  \item The GRU candidate state was written as
        $\tanh(W_n[x_t + r_t\odot h_{t-1}] + b_n)$. The correct form is
        $\tanh(W_{in}x_t + b_{in} + r_t\odot(W_{hn}h_{t-1} + b_{hn}))$.
  \item The text claimed one spline per input dimension; a KAN layer has one
        learnable function per edge, i.e.\ $d_{in}\times d_{out}$ functions.
  \item ``Parameter efficiency'' was listed as a KAN advantage, and the title
        used ``Lightweight'', although both KAN models were larger and slower
        than the baseline.
  \item The reference to Hou et al.\ described a survey; the cited version of
        arXiv:2407.11075 is a critical assessment (by Hou, Ji, Zhang and
        Stefanidis) reporting that KANs are substantially slower than MLPs and
        rarely more accurate outside symbolic regression. Rendle et al.\ (2010)
        was listed in the bibliography but never cited.
  \item The research gap claimed no prior work on KAN-based recurrent cells;
        TKAN (Genet and Inzirillo, 2024) is such a cell. The claim has been
        narrowed (Chapter~2).
\end{enumerate}
